\documentclass{article}
\usepackage{amsmath,amssymb,graphicx,booktabs,hyperref,subcaption,float}
\graphicspath{{figures/}}
\begin{document}

\title{Supporting Information: Cooperative Yaw Control for Wind Farm Power Maximization}
\maketitle

\section{SAC Algorithm Comparison --- Full Details}

\subsection{Implementation}
A JAX/NNX-native Soft Actor-Critic (SAC) agent was implemented on the same on-device stack used for PPO. The SAC agent uses twin Q-networks with $[128,128]$ MLPs (identical architecture to PPO), a Gaussian policy with tanh squashing, automatic entropy tuning ($\alpha_0{=}0.2$, target entropy ${=}{-}\mathrm{dim}(\mathcal{A})$), and a replay buffer of $5{\times}10^{5}$ transitions. Training used $N_\mathrm{envs}{=}4$ parallel environments (vs.\ PPO's 128; the lower count is a result of JIT compilation constraints for SAC's multi-gradient update structure) and $N_\mathrm{steps}{=}512$ steps per rollout. Five seeds were trained for $3{\times}10^{7}$ environment steps each under Config-E settings (J=3, deficit normalization, position encoding, SLSQP-regret reward, focused wind sampling 0.3/0.3/0.4, $\pm10^\circ$ action bounds, downstream locking).

\subsection{Results}
At the matched $3{\times}10^{7}$-step budget, SAC achieves a five-seed-mean final episode return of $+185.2\,\pm\,75.3$ (regret-reward units), compared to PPO Config-E's $+334.3\,\pm\,217.0$. SAC reaches $55.4\%$ of PPO's mean reward with substantially lower inter-seed variance ($\pm75.3$ vs.\ $\pm217.0$), indicating more consistent but overall lower performance. The entropy coefficient $\alpha$ converges from $0.2$ to $0.010$--$0.028$ across seeds. The throughput gap is $13{\times}$ (${\sim}7{,}100$\,FPS for SAC vs.\ ${\sim}95{,}000$\,FPS for PPO), arising from SAC's three separate gradient computations per update.

\subsection{Discussion}
We attribute SAC's underperformance to: (i)~off-policy sample inefficiency when the reward signal is concentrated in a narrow wind regime (${\sim}9\%$ of conditions); (ii)~the squashed-Gaussian action parameterization penalizing large yaw increments; and (iii)~automatic entropy tuning requiring substantial interaction data before exploitation. A pure-JAX refactoring of the SAC gradient computation (bypassing NNX state merging inside traced functions, reducing per-step overhead from $56$\,ms to $14$\,ms) is expected to narrow the throughput gap.

\section{Behavioral Cloning Baseline --- Full Details}

\subsection{Implementation}
A Behavioral Cloning (BC) baseline was trained via supervised regression to predict SLSQP-optimal yaw vectors from wind conditions, using a dataset of 1{,}000 SLSQP solutions sampled from the training distribution. The BC policy uses the same $[128,128]$ MLP architecture as the PPO actor and was trained for 100 epochs with MSE loss on the yaw-angle targets.

\subsection{Results}
Evaluated on 100 uniformly sampled wind conditions, the BC policy achieves a marginal mean gain of $\mathbf{-20.20\%}$ and an aligned-cube gain of $-31.50\%$ ($n{=}7$ aligned-cube conditions)---substantially worse than the zero-yaw baseline. The negative marginal gain indicates that the BC policy actively reduces farm power relative to doing nothing.

\subsection{Discussion}
The BC result suggests that direct supervised regression to SLSQP yaw vectors is insufficient under the present dataset ($1{,}000$ samples) and architecture ($[128,128]$ MLP), likely because the SLSQP yaw mapping is discontinuous near regime boundaries where turbines enter or exit the wake. A feedforward MLP trained on a finite sample may not capture these discontinuities, whereas the DRL policy learns through interaction and produces smoother yaw trajectories. Further work with larger expert datasets and specialized architectures would be needed to draw stronger conclusions about the relative merits of BC versus DRL for this problem.

\section{JAX/NNX Implementation Efficiency}

The initial Stable-Baselines3 PPO training loop was profiled, revealing that environment stepping dominates wall-clock time (${\sim}86\%$) while the PPO policy gradient accounts for only ${\sim}14\%$. This motivated a full on-device rewrite: the wake model was re-implemented in pure JAX, vectorized over turbines and parallel environments with \texttt{vmap}, and the entire (rollout + GAE + PPO update) loop was JIT-compiled. The result is a ${\sim}30{\times}$ throughput jump to ${\sim}10^{5}$\,FPS at $N_\mathrm{envs}{=}128$. A five-seed $6{\times}10^{7}$-step PPO run completes in ${\sim}10$\,min per seed on an NVIDIA RTX 4090. SAC training achieves ${\sim}7{,}100$\,FPS at $N_\mathrm{envs}{=}4$.

\section{KL Early-Stopping Ablation}

To quantify the independent contribution of KL-targeted early stopping, two Config-D variants were trained for $3{\times}10^{7}$ steps (3 seeds each): \textsc{kl\_on} (TARGET\_KL=0.015, default) and \textsc{kl\_off} (TARGET\_KL=100, effectively disabled). All other hyperparameters match Config-D. Table~S1 summarizes the final training metrics extracted from the saved checkpoints.

\begin{table}[h]
\centering
\caption{KL early-stop ablation: final training metrics (3-seed mean $\pm$ std).}
\label{tab:kl_ablation}
\begin{tabular}{lcc}
\toprule
Metric & KL ON (0.015) & KL OFF (100) \\
\midrule
Final episode reward & $+344.0\,\pm\,93.7$ & $+311.0\,\pm\,138.6$ \\
Mean KL divergence (last iter.) & $0.0021$ & $0.0021$ \\
Early-stop trigger rate & $0.3\%$ & $0.0\%$ \\
\bottomrule
\end{tabular}
\end{table}

The KL ON configuration achieves a modestly higher final episode reward ($+344.0$ vs.\ $+311.0$) with slightly lower inter-seed variance. The early-stop mechanism triggers on only $0.3\%$ of iterations under the default TARGET\_KL=0.015, confirming that it functions as a safety valve rather than a primary regularization mechanism under Config-D. The near-zero trigger rate explains why the mean KL divergence is identical between the two configurations: the KL threshold is rarely exceeded, so the training trajectories are nearly indistinguishable. Cosine learning-rate decay and AdamW weight decay are therefore the dominant contributors to Config-D's improved training stability, with KL early-stopping providing a marginal additional safeguard.

\section{Supplementary Figures}
\begin{figure}[t]
\centering
\includegraphics[width=0.85\linewidth]{fig_methods_comparison}
\caption{Schematic comparison of representative wind-farm cooperative yaw control paradigms.}
\label{fig:methods}
\end{figure}

\begin{figure}[t]
\centering
\includegraphics[width=0.95\linewidth]{fig_lock_ablation_eval}
\caption{Lock-mask ablation on the $3\!\times\!3$ layout, three seeds (Config-A baseline configuration). Mean gain over uniformly sampled wind conditions: $\mathbf{+2.42\%}$ (\textsc{lock\_on}) vs.\ $\mathbf{-0.46\%}$ (\textsc{lock\_off}), with negative-gain fractions of $18\%$ and $28\%$ respectively. Locking is structurally necessary for stable cooperative-yaw learning. The lock mechanism is retained in all subsequent configurations including Config-E.}
\label{fig:lock_ablation}
\end{figure}

\begin{figure}[t]
\centering
\includegraphics[width=0.95\linewidth]{fig_nnx_vs_sb3}
\caption{Wall-clock A/B benchmark on the $3\!\times\!3$ PPO training loop. SB3-PyTorch on CPU/CUDA, NNX-JAX on CPU/CUDA with the NumPy environment, and the fully on-device NNX+JAX-Env stack are compared at matched total environment steps. Throughput jumps by $\sim$$30\!\times$ only once the environment itself is ported to JAX --- swapping the policy backend alone is throughput-neutral because the NumPy environment occupies $86\%$ of wall-clock and the host$\leftrightarrow$device boundary cancels the GPU's compute savings.}
\label{fig:nnx_vs_sb3}
\end{figure}

\begin{figure}[t]
\centering
\includegraphics[width=0.95\linewidth]{fig_xval_jaxenv_policy_gain}
\caption{Numerical cross-validation of the on-device JAX environment against the original NumPy gym. A single short-budget ($5\!\times\!10^{6}$ steps) reference policy is evaluated on a $19$-point $(\phi,v)$ grid inside both environments. The two implementations agree on sign, rank ordering and the location of the peak gain, supporting the use of the on-device stack for accelerated training without loss of physical fidelity.}
\label{fig:xval_env}
\end{figure}

\begin{figure}[h]
\centering
\begin{subfigure}{0.32\linewidth}\includegraphics[width=\linewidth]{fig_3turbine_layout}\caption{Three-turbine inline.}\label{fig:3t}\end{subfigure}\hfill
\begin{subfigure}{0.32\linewidth}\includegraphics[width=\linewidth]{fig_large_farm_layout}\caption{Large 35-turbine farm.}\label{fig:largefarm}\end{subfigure}\hfill
\begin{subfigure}{0.32\linewidth}\includegraphics[width=\linewidth]{fig_2turbine_layout}\caption{Two-turbine yaw.}\label{fig:2t}\end{subfigure}
\caption{Calibration CFD reference datasets.}
\end{figure}

\begin{figure}[t]
\centering
\begin{subfigure}{0.48\linewidth}\includegraphics[width=\linewidth]{fig_hornsrev_layout}\caption{Horns Rev layout.}\label{fig:hr_layout}\end{subfigure}\hfill
\begin{subfigure}{0.48\linewidth}\includegraphics[width=\linewidth]{fig_v80_power_curve}\caption{Predicted V-80 power curve.}\label{fig:v80}\end{subfigure}
\caption{Horns Rev calibration assets.}
\end{figure}

\begin{figure}[t]
\centering
\begin{subfigure}{0.49\linewidth}\includegraphics[width=\linewidth]{fig_floris_hornsrev_compare}\caption{$3\!\times\!3$ NREL-5MW direction sweep, $\gamma\!=\!0$, $U_\infty\!=\!11.4\,\mathrm{m/s}$.}\label{fig:floris_sweep}\end{subfigure}\hfill
\begin{subfigure}{0.49\linewidth}\includegraphics[width=\linewidth]{fig_floris_yaw_sweep}\caption{Two-turbine yaw sweep, $\gamma\!\in\![0^\circ,50^\circ]$.}\label{fig:floris_yaw}\end{subfigure}
\caption{Cross-validation of the proposed gray-box wake model against FLORIS. (a) Direction-by-direction farm power on the $3\!\times\!3$ NREL-5MW grid. (b) Two-turbine yaw sweep capturing the yaw-deflection regime exploited by the closed-loop controller.}
\label{fig:floris_xval}
\end{figure}

\begin{figure}[t]
\centering
\begin{subfigure}{0.48\linewidth}\includegraphics[width=\linewidth]{fig_1x2_training_curve}\caption{Training curve.}\label{fig:1x2_train}\end{subfigure}\hfill
\begin{subfigure}{0.48\linewidth}\includegraphics[width=\linewidth]{fig_1x2_optimization_rate}\caption{Wind-rose optimization rate.}\label{fig:1x2_opt}\end{subfigure}
\caption{$1{\times}2$ inline farm: PPO training and full wind-rose performance.}
\end{figure}

\begin{figure}[t]
\centering
\includegraphics[width=0.7\linewidth]{fig_3x3_optimization_rate}
\caption{Full wind-rose optimization rate for the $3{\times}3$ layout, LCB-selected PPO policy.}
\label{fig:3x3_opt}
\end{figure}

\begin{figure}[t]
\centering
\includegraphics[width=0.95\linewidth]{fig_yaw_rose}
\caption{Learned yaw angles across the wind rose for the $3{\times}3$ layout (seed~0), shown at $v\!=\!8\,\mathrm{m/s}$ (top) and $v\!=\!11.4\,\mathrm{m/s}$ (bottom). Cooperative yaw deflection is active only within the column-aligned band $|\phi\!-\!270^\circ|\!<\!15^\circ$ (shaded); outside this band the policy correctly reverts to near-zero yaw. Crosses mark turbines identified as downstream by the locking mechanism.}
\label{fig:yaw_rose}
\end{figure}

\begin{figure}[t]
\centering
\includegraphics[width=0.95\linewidth]{fig_drl_vs_slsqp_yaw_comparison}
\caption{Per-turbine yaw angle comparison between the DRL policy (\textsc{Config-E}, seed~1) and the SLSQP optimum across the aligned-cube wind regime ($\phi\!\in\![255^\circ,285^\circ]$, $v\!\in\!\{7,9,11\}\,\mathrm{m/s}$). The DRL policy closely matches SLSQP at off-axis directions (mean $|\Delta|\!=\!1.6^\circ$) but is systematically more conservative at the perfectly-aligned condition $\phi\!=\!270^\circ$ (mean $|\Delta|\!=\!17.1^\circ$), explaining the residual $0.26$\,pp recovery gap.}
\label{fig:drl_vs_slsqp_yaw}
\end{figure}

\begin{figure}[t]
\centering
\includegraphics[width=0.95\linewidth]{fig_inference_latency}
\caption{Measured single-CPU-core inference latency (MLP forward pass only, excluding sensor/communication/actuator latencies) of the trained PyTorch policy over $3000$ random observations per layout. Median (p50) latency is $0.114$/$0.116$/$0.172\,\mathrm{ms}$ for $N\!=\!2$/$9$/$25$; p95 stays below $0.20\,\mathrm{ms}$ across all layouts. Latency is approximately flat in $N$ within the hidden-width regime used here.}
\label{fig:infer_lat}
\end{figure}

\begin{figure}[t]
\centering
\includegraphics[width=\linewidth]{figures/fig_reward_penalty_ablation.pdf}
\caption{\textbf{Reward-penalty ablation on the $3{\times}3$ layout.}
Single-panel episode-return learning curves (mean $\pm$ min/max envelope over $3$ seeds,
$3\!\times\!10^{7}$ environment steps each).
Adding $P_\mathrm{mag}+P_\mathrm{rate}$ with
$\lambda_\mathrm{mag}=\lambda_\mathrm{rate}=8\!\times\!10^{-5}$
depresses the final-$10$-iteration return from $+19.79$ to $+15.78$
($-20\%$ relative) and lowers the best-seed peak from $+28.1$ to $+22.3$.
Per-seed final returns and per-group entropy diagnostics are saved in
\texttt{reward\_penalty\_ablation.json}.}
\label{fig:reward_ablation}
\end{figure}

\begin{figure}[t]
\centering
\includegraphics[width=0.95\linewidth]{fig_obs_noise_robustness}
\caption{Observation noise robustness. (a)~Gain vs.\ wind-direction noise $\sigma_\phi$. (b)~Gain vs.\ wind-speed noise $\sigma_v$. (c)~2D heatmap of aligned-cube gain as a function of $(\sigma_\phi,\sigma_v)$. The controller tolerates $\sigma_\phi\!=\!10^\circ$ with $<\!1\%$ relative gain erosion.}
\label{fig:noise_robustness}
\end{figure}

\section{J=15 Extended Observation Horizon --- Causal Ablation and J-Trend}

\subsection{J-horizon trend under dynamic wind}
Table~S2 summarizes the dynamic-wind performance across observation history lengths $J\in\{3,8,15\}$ under both AR(1) protocols (Proto~A: $\rho{=}0.99$, $\sigma{=}1^\circ$; Proto~B: $\rho{=}0.95$, $\sigma{=}2^\circ$). All policies use $\lambda_\mathrm{rate}{=}5{\times}10^{-4}$ and $6{\times}10^{7}$ training steps (5 seeds each). The mean per-step gain improves monotonically with $J$: from $-1.52\%$ (J=3, Proto~A) to $-0.94\%$ (J=15, Proto~A), and from $-1.21\%$ to $-0.53\%$ under Proto~B. Yaw travel also decreases with longer observation windows ($-27\%$ from J=3 to J=15), consistent with better temporal context reducing reactive over-steering. The J=15 Proto~B result ($-0.53\%$) is $2.9{\times}$ better than J=3 Proto~B ($-1.21\%$).

\begin{table}[h]
\centering
\caption{Dynamic-wind performance vs.\ observation history length $J$ (5-seed mean, $\lambda_\mathrm{rate}{=}5{\times}10^{-4}$, $6{\times}10^{7}$ steps).}
\label{tab:j_trend}
\begin{tabular}{lcccc}
\toprule
Configuration & $J$ & Proto~A gain [\%] & Proto~B gain [\%] & Travel/turb [$^\circ$] \\
\midrule
Static J=3 (reference) & 3 & $-0.178$ & $-0.120$ & $180$ \\
Dynamic $\lambda{=}5{\times}10^{-4}$ & 3 & $-1.519$ & $-1.210$ & $376$ \\
Dynamic $\lambda{=}5{\times}10^{-4}$ & 8 & $-1.667$ & $-1.231$ & $344$ \\
Dynamic $\lambda{=}5{\times}10^{-4}$ & 15 & $-0.937$ & $-0.534$ & $273$ \\
\bottomrule
\end{tabular}
\end{table}

\subsection{Causal ablation of temporal observation ordering}
The causal ablation experiment reported in the main text (Section~4.8, Table~16) confirms that the J=15 policy exploits temporally-ordered observation history rather than merely benefiting from a larger input dimension. The complete results are available in \texttt{results/j15\_causal\_ablation.json}.

\section{Load Proxy Analysis}

To provide first-order structural-load context without full aeroelastic simulation, we compute two simplified yaw-activity load proxies on the unified dynamic-wind benchmark (200 AR(1) trajectories, 200 steps, 5 Config-E seeds). Proxy~B (cumulative squared yaw change, $\sum\Delta\gamma^2$) approximates yaw-bearing wear proportional to the squared displacement per step; Proxy~C (cumulative absolute yaw rate, $\sum|\Delta\gamma|/\Delta t$) approximates gearbox fatigue proportional to the integrated yaw speed. Table~S3 compares Static Config-E, Dynamic $\lambda{=}5{\times}10^{-4}$, and the zero-yaw baseline.

\begin{table}[h]
\centering
\caption{Yaw-activity load proxies on the unified dynamic-wind benchmark (5-seed Config-E mean).}
\label{tab:load_proxy}
\begin{tabular}{lccc}
\toprule
Controller & Proxy~B $\sum\Delta\gamma^2$ & Proxy~C $\sum|\Delta\gamma|/\Delta t$ & Reversals \\
\midrule
Zero yaw (baseline) & $0$ & $0$ & $0$ \\
Static Config-E & $2.54{\times}10^{6}$ & $4{,}007$ & $1{,}662$ \\
Dynamic $\lambda{=}5{\times}10^{-4}$ & $1.35{\times}10^{7}$ & $10{,}419$ & $1{,}684$ \\
\bottomrule
\end{tabular}
\end{table}

Static training followed by post-hoc DRL-Deploy gating achieves $5.3{\times}$ lower Proxy~B and $2.6{\times}$ lower Proxy~C compared to dynamic-wind training at $\lambda_\mathrm{rate}{=}5{\times}10^{-4}$, suggesting a more favorable power--load trade-off for the static+gate strategy. These proxies are first-order approximations; full aeroelastic validation using OpenFAST or FAST.Farm is required for definitive structural-load conclusions.

\end{document}
